% Included by project_log.tex. Regenerate with
%   ./venv/bin/python src/build_dataset.py && ./venv/bin/python src/run_triage.py

\section{Phase 4: risk triage classifier}\label{sec:phase4}

\subsection{Manufacturing the negative class}

\texttt{cdm\_public} contains no negative examples: every row is
\texttt{EMERGENCY\_REPORTABLE} with $\texttt{PC} > 10^{-4}$. A classifier
trained on it learns to answer ``yes'' and scores \SI{100}{\percent}.
Negatives were therefore manufactured by screening debris catalogues
directly -- dense, co-orbital populations that generate many close
approaches per object.

\begin{table}[htbp]
\centering
\caption{Screened dataset: three debris families, seven days each, 60\,s
         coarse grid, refined to \SI{50}{\kilo\meter}.}
\label{tab:dataset}
\begin{tabular}{lrrr}
\toprule
Catalogue & Objects & Conjunctions & $\log_{10}\Pc > -10$ \\
\midrule
Fengyun-1C debris   & 1971 & 308{,}264 & 324 \\
Cosmos-2251 debris  &  584 &  28{,}255 &  26 \\
Iridium-33 debris   &  111 &   1{,}270 &   2 \\
\midrule
Total               & 2666 & 337{,}789 & 352 \\
\bottomrule
\end{tabular}
\end{table}

The positive rate is \SI{0.104}{\percent}. $\log_{10}\Pc$ spans
$-29843$ to $-4.5$, with a median of $-5422$: most screened pairs miss by
tens of standard deviations, where $\Pc$ underflows to exactly zero in
double precision. Labels are therefore computed directly in logarithms,
without ever forming $\Pc$, or the entire negative class collapses to a
single value.

\gotcha{\textbf{Not one conjunction reached the operational
$\Pc > 10^{-4}$ threshold.} A debris cloud over a one-week window simply
contains no operationally red events, so the label threshold must be
justified rather than inherited. Results below use $10^{-10}$; the ordering
of methods is unchanged at $10^{-8}$ and $10^{-12}$.}

\subsection{Screening cost and the linear-motion filter}

The coarse gate must be wide. At \SI{16}{\kilo\meter\per\second} closing
speed a \SI{60}{\second} grid steps \SI{960}{\kilo\meter}, so a pair passing
within \SI{1}{\kilo\meter} may never appear closer than
$\sim$\SI{450}{\kilo\meter} at any sampled instant. Screening at a narrow
threshold on a coarse grid therefore misses almost every genuine approach.
The gate used is
\begin{equation}
  d_{\text{gate}} = d_{\text{target}}
    + \tfrac{1}{2}\, v_{\max}\, \Delta t,
  \label{eq:gate}
\end{equation}
which is \SI{530}{\kilo\meter} for the parameters above -- and which leaves
essentially every pair a candidate.

\worked{An analytic filter makes that affordable. Near closest approach
relative motion is nearly rectilinear, so from sampled relative position
$\Delta r$ and velocity $\Delta v$ the minimum follows in closed form at
$t^{*} = -(\Delta r \cdot \Delta v)/\|\Delta v\|^2$. Applying this to every
local minimum removed \SI{91}{\percent} of candidates (163{,}806 to
14{,}528) with \emph{byte-identical} output, cutting a Cosmos-2251 day from
\SI{29.5}{\second} to \SI{7.2}{\second}.}

\subsection{Evaluation design}

Accuracy and AUC are near-meaningless at a \SI{0.1}{\percent} positive rate;
a constant ``no'' scores \SI{99.9}{\percent}. The operational question is
asymmetric -- missing a high-risk conjunction is unacceptable, passing a
harmless one costs only compute -- so models are compared by \textbf{how
much of the catalogue they discard at full recall}.

Two design choices guard against self-deception:
\begin{enumerate}
  \item \textbf{The baseline is a miss-distance cut}, not random. That is
        what an analyst does without a model. A classifier that cannot beat
        one threshold on one feature is not worth having.
  \item \textbf{Splits are grouped, not random.} A day split holds out the
        last two days; an object split assigns objects to train or test and
        keeps only pairs whose \emph{both} members are held out, discarding
        cross pairs. Random splitting would leak the same object pair into
        both halves at slightly different times.
\end{enumerate}

\begin{table}[htbp]
\centering
\caption{Triage performance. ``Kept'' is the fraction of conjunctions passed
         on for expensive analysis; lower is better at equal recall.
         Thresholds are set from out-of-fold predictions on training data,
         then applied unchanged.}
\label{tab:triage}
\begin{tabular}{llrrr}
\toprule
Split & Model & Kept & Recall & Missed \\
\midrule
\multirow{3}{*}{Day}
 & Miss-distance cut (baseline) & \SI{1.41}{\percent} & \SI{100.0}{\percent} & 0 \\
 & Logistic regression          & \SI{0.92}{\percent} & \SI{100.0}{\percent} & 0 \\
 & Gradient boosting            & \SI{100.0}{\percent} & \SI{100.0}{\percent} & 0 \\
\midrule
\multirow{3}{*}{Object}
 & Miss-distance cut (baseline) & \SI{1.15}{\percent} & \SI{97.2}{\percent}  & 1 \\
 & Logistic regression          & \SI{0.82}{\percent} & \SI{100.0}{\percent} & 0 \\
 & Gradient boosting            & \SI{100.0}{\percent} & \SI{100.0}{\percent} & 0 \\
\bottomrule
\end{tabular}
\end{table}

\subsection{Reading the result}

\worked{Logistic regression beats the analyst's baseline on both splits. It
passes \SIrange{0.82}{0.92}{\percent} of conjunctions for expensive analysis
against \SIrange{1.15}{1.41}{\percent} for a miss-distance cut -- roughly
\textbf{a third fewer expensive analyses at equal or better recall}.}

\worked{More importantly, on the object split the \textbf{baseline fails
full recall} (\SI{97.2}{\percent}, one high-risk conjunction missed) while
logistic regression retains all of them. Miss distance alone is not
sufficient: its rank correlation with $\log_{10}\Pc$ is only $+0.60$,
because the projected covariance depends on encounter geometry that miss
distance does not capture.}

\failed{\textbf{Gradient boosting fails completely}, retaining
\SI{100}{\percent} of the catalogue under every configuration. This is not a
threshold artefact: its \emph{oracle} operating point is also
$\approx$\SI{100}{\percent}, meaning the lowest-scoring true positive is
ranked below almost the entire test set. The flexible model fragments a
space in which the rare positives live in a thin tail, and cannot rank them;
the linear model, whose score is monotone in roughly the right direction,
can. At \SI{0.1}{\percent} positive rate, model capacity actively hurts.}

\gotcha{An earlier version of this evaluation set thresholds from
\emph{in-sample} training scores. Gradient boosting fits its training
positives easily, so the implied threshold did not transfer and the
operating point collapsed. Any threshold intended for deployment must be
chosen out-of-fold. This also masked the baseline's recall failure.}

\subsection{Limitations}

\begin{enumerate}
  \item Test sets contain only 36--99 positives. Differences of a few tenths
        of a percent in ``kept'' should not be over-read.
  \item The population is debris clouds, not a general catalogue. Debris is
        small, uncontrolled, and co-orbital; operational payload
        conjunctions differ.
  \item Labels inherit every Phase 2 assumption -- the synthesized
        covariance and a fixed \SI{2}{\meter} hard-body radius. They are
        self-consistent, not ground truth.
  \item Seven-day propagation from a single epoch exceeds the
        \numrange{2}{5}~day range over which the covariance model was
        calibrated.
\end{enumerate}
